%% POLISHED ABSTRACT & INTRODUCTION FOR ICLR 2027
%% Copy-paste these sections into main.tex
%% Author: Ege Berk Türk
%% Date: May 23, 2026

%%%%%%%%%%%%%%%%%%%%%%%%%%%%%%%%%%%%%%%%%%%%%%%%%%%%%%%%%%%%%%%%
%% ABSTRACT (POLISHED VERSION)
%%%%%%%%%%%%%%%%%%%%%%%%%%%%%%%%%%%%%%%%%%%%%%%%%%%%%%%%%%%%%%%%

\begin{abstract}
High-dimensional embeddings ($d \geq 1536$) dominate modern NLP and retrieval systems, yet their storage costs remain prohibitive. We introduce \textbf{Phase-Multiplexed Spherical Collapse (PMSC)}, a training-free lossless compression method that achieves \textbf{3.50$\times$ compression} by exploiting three key insights: (1) \textit{angular concentration} — unit-norm vectors' spherical coordinates cluster near $\pi/2$, causing IEEE 754 exponent uniformity; (2) \textit{delta transformation} — subtracting this concentration center ($\varphi - \pi/2$) exposes mantissa-level redundancy; (3) \textit{FFT basis sharing} — 2D reshaping and frequency-domain processing enable a shared amplitude basis across vector batches, storing only per-vector phase keys. 

On 10,000 random $d=1536$ vectors, PMSC achieves 3.50$\times$ bit-perfect compression (vs. 1.50$\times$ baseline), $\text{exponent\_127\_pct} = 99.75\%$, and delta entropy of 3.15 bits/element—a 10$\times$ reduction. Real semantic embeddings (MiniLM, MPNet) exhibit identical behavior. In lossy float16 mode, cosine similarity remains 1.0000 with 2.09$\times$ compression. PMSC decodes in $<$1ms/vector, making it production-ready for RAG systems and vector databases.
\end{abstract}

%%%%%%%%%%%%%%%%%%%%%%%%%%%%%%%%%%%%%%%%%%%%%%%%%%%%%%%%%%%%%%%%
%% INTRODUCTION (NEW SECTION)
%%%%%%%%%%%%%%%%%%%%%%%%%%%%%%%%%%%%%%%%%%%%%%%%%%%%%%%%%%%%%%%%

\section{Introduction}
\label{sec:introduction}

Embedding vectors power modern machine learning—from semantic search~\cite{reimers2019sentencebert} to retrieval-augmented generation (RAG)~\cite{lewis2020retrieval}. However, storing billions of high-dimensional vectors ($d \geq 1536$) at 6KB each (float32) is economically unsustainable. While product quantization (PQ)~\cite{jegou2010product} achieves $\sim$6$\times$ compression, it introduces irreversible quantization errors, degrading retrieval quality.

Recent work by Xiao~\cite{xiao2026spherical} identified ``spherical collapse''—a phenomenon where angular coordinates of unit-norm vectors concentrate near $\pi/2$, causing IEEE 754 exponent uniformity. By storing only exponent-127 values, their method achieves \textbf{1.50$\times$ lossless compression}. Yet this only scratches the surface of the underlying structure.

\subsection{Our Contributions}

We present \textbf{Phase-Multiplexed Spherical Collapse (PMSC)}, which improves upon the baseline in three ways:

\begin{enumerate}
\item \textbf{Delta Transformation:} Subtracting the concentration center ($\delta_k = \varphi_k - \pi/2$) shifts values to near-zero, exposing mantissa-level redundancy and reducing entropy from 32 bits to 3.15 bits/element.
\item \textbf{FFT Basis Sharing:} Reshaping delta vectors into 2D matrices and applying 2D FFT reveals frequency-domain patterns. A \textit{shared amplitude basis} across batches enables storing only per-vector phase keys, eliminating inter-vector redundancy.
\item \textbf{Adaptive Optimization:} Automatic keep\_ratio tuning based on delta entropy maximizes compression without hyperparameter search.
\end{enumerate}

\textbf{Results:} On 10K random $d=1536$ vectors, PMSC achieves:
\begin{itemize}
\item \textbf{3.50$\times$ bit-perfect compression} (2.33$\times$ better than baseline)
\item Sub-millisecond decoding (0.77ms/vector)
\item Validated on real semantic embeddings (MiniLM, MPNet): 98-99\% exponent\_127\_pct
\item Lossy float16 mode: 2.09$\times$ compression with 1.0000 cosine similarity
\end{itemize}

Unlike PQ or scalar quantization (SQ), PMSC is \textit{training-free}, \textit{deterministic}, and \textit{bit-perfect}. Our code is open-source with 20+ passing tests.

\subsection{Paper Organization}

Section~\ref{sec:related} reviews prior work. Section~\ref{sec:theory} formalizes spherical collapse and proves compression bounds. Section~\ref{sec:method} details PMSC's pipeline. Section~\ref{sec:experiments} presents 10K-scale benchmarks. Section~\ref{sec:ablation} quantifies component contributions. Section~\ref{sec:conclusion} discusses future work.

%%%%%%%%%%%%%%%%%%%%%%%%%%%%%%%%%%%%%%%%%%%%%%%%%%%%%%%%%%%%%%%%
%% USAGE INSTRUCTIONS
%%%%%%%%%%%%%%%%%%%%%%%%%%%%%%%%%%%%%%%%%%%%%%%%%%%%%%%%%%%%%%%%
%% 
%% 1. Replace the existing abstract in main.tex (lines 21-22) with the
%%    POLISHED ABSTRACT above.
%%
%% 2. Insert the INTRODUCTION section after \maketitle and before 
%%    the first existing section (probably Related Work).
%%
%% 3. Update bibliography:
%%    - Add \cite{reimers2019sentencebert} (Sentence-BERT paper)
%%    - Add \cite{lewis2020retrieval} (RAG paper)
%%    - Verify \cite{jegou2010product} (FAISS PQ)
%%    - Verify \cite{xiao2026spherical} (arXiv 2602.00079)
%%
%% 4. Ensure section labels are consistent:
%%    - \label{sec:related}
%%    - \label{sec:theory}
%%    - \label{sec:method}
%%    - \label{sec:experiments}
%%    - \label{sec:ablation}
%%    - \label{sec:conclusion}
%%
%%%%%%%%%%%%%%%%%%%%%%%%%%%%%%%%%%%%%%%%%%%%%%%%%%%%%%%%%%%%%%%%
